\documentclass[11pt, a4paper]{article}

\usepackage[top=2cm, bottom=2cm, left=2.5cm, right=2cm]{geometry}
\usepackage{fontspec}
\usepackage{amsmath, amssymb}
\usepackage{booktabs}
\usepackage{tabularx}
\usepackage{graphicx}
\usepackage{float}
\usepackage{xcolor}
\usepackage{hyperref}
\usepackage{xurl}
\usepackage{parskip}
\usepackage{listings}
\usepackage{microtype}
\usepackage{caption}
\usepackage{enumitem}

\graphicspath{{images/}}
\emergencystretch=3em

\hypersetup{
    colorlinks=true,
    linkcolor=blue!70!black,
    citecolor=blue!70!black,
    urlcolor=blue!70!black,
    pdftitle={Assignment 05 - CNN-3 vs CNN-5 in TensorFlow and PyTorch},
    pdfauthor={Nguyen Van Truong}
}

\lstset{
    basicstyle=\small\ttfamily,
    breaklines=true,
    frame=single,
    backgroundcolor=\color{gray!10},
    columns=fullflexible
}

\newcommand{\code}[1]{\texttt{\detokenize{#1}}}
\newcommand{\run}[1]{\texttt{#1}}
\setlist{itemsep=2pt, topsep=4pt}

\begin{document}

% ---------------------------------------------------------
% COVER PAGE
% ---------------------------------------------------------
\begin{titlepage}
\begin{center}

\vspace*{0.5cm}

\Large
\textbf{POSTS AND TELECOMMUNICATIONS INSTITUTE OF TECHNOLOGY}\\
\textbf{Faculty of Information Technology}

\vspace{2.0cm}

\Huge
\textbf{Intelligent System Development}\\
\vspace{0.4cm}
\LARGE
\textbf{ASSIGNMENT 05}\\
\vspace{0.6cm}
\Large
\textbf{Convolutional Neural Networks:\\CNN-3 vs.\ CNN-5 in TensorFlow/Keras and PyTorch}

\vspace{1.2cm}

\vfill

\normalsize
\begin{tabular}{rl}
\textbf{Instructor / Lecturer:} & Assoc. Prof. Dinh Que Tran, Ph.D. \\
\textbf{Email:} & quetd@ptit.edu.vn, tdque@yahoo.com \\[0.8cm]
\textbf{Student Name:} & Nguyễn Văn Trường \\
\textbf{Student ID:} & B23DCCE095 \\
\textbf{Semester:} & I.2025 -- 2026 \\
\textbf{Class:} & E23CNPM02 \\
\textbf{Submission:} & Individual Project \\
\end{tabular}

\vfill

\textbf{Hanoi, 2026}

\vspace{0.5cm}

\end{center}
\end{titlepage}

% ---------------------------------------------------------
% TABLE OF CONTENTS
% ---------------------------------------------------------
\tableofcontents
\newpage

% =====================================================================
\section{Executive Summary}
% =====================================================================

This project sets up convolutional networks of \textbf{3} and \textbf{5} conv layers
(\run{CNN-3}, \run{CNN-5}) using both \textbf{TensorFlow/Keras} (\run{TF}) and \textbf{PyTorch} (\run{PT}), then fits
them to three datasets: \textbf{SVHN} (99,289 house-number photos, 10 classes), \textbf{GTSRB} (51,839 traffic-sign
photos, 43 classes) and a \textbf{diabetes} table (300,000 patients, 24 features, binary target). Every matched pair of
runs shares the split, architecture, parameter count, optimizer, batch size, early-stopping rule and seed, so the only
things that change are the depth of the network and the framework used. Each trained model is tagged with a run ID of
the form \run{<dataset>-<framework>-<model>}, for instance \run{SV-PT-CNN-5}.

\begin{table}[H]
\centering
\caption{All 12 runs (test set). s/epoch = median training time per epoch, CPU.}
\label{tab:summary}
\small
\input{summary_table.tex}
\end{table}

\textbf{Key takeaways}
\begin{enumerate}
\item On the two image datasets, adding depth helps: CNN-5 beats CNN-3 under both frameworks, with macro F1 up by
  +0.021 to +0.040 on SVHN and +0.015 to +0.036 on GTSRB. The best model overall is \run{GT-PT-CNN-5}, with a
  macro F1 of 0.9979.
\item The diabetes dataset tells a different story. All four CNN variants sit in a narrow band of macro F1
  (0.627--0.630) and ROC-AUC (0.694--0.696), and a much smaller MLP (6,977 parameters versus 391,329) reaches
  ROC-AUC 0.695 as well. The features themselves carry a weak signal, so no amount of extra depth changes that.
\item Framework choice has a modest, dataset-dependent effect: TF and PT macro F1 stay within 0.013 of each other on
  every dataset, though which framework comes out ahead flips between CNN-3 and CNN-5, and their parameter counts are
  identical either way (checked directly in code).
\item The cost of depth is time, not accuracy loss: CNN-5 needs 1.8--2.2$\times$ as long per epoch as CNN-3. On this
  CPU-only run, PyTorch was 37--43\% slower per epoch than Keras on the two image datasets and roughly 2$\times$ slower
  on the tabular one.
\end{enumerate}

% =====================================================================
\section{CNN Fundamentals}
% =====================================================================

Each term below is defined a single time and keeps that meaning throughout the rest of the report.
Notation used: $N$ samples; input $C\times H\times W$ (channels, height, width); kernel size $K$, stride $S$,
padding $P$; $F$ filters (output channels); logits $z$, predicted probability $p$, true label $y$.

\subsection{Convolution}
A convolutional layer passes $F$ learnable \textbf{filters}, each of size $K\times K\times C_{in}$, across the input.
At every position it produces a single value belonging to a \textbf{feature map}:
\[
\text{out}(i,j) = \sum_{u,v,c} \text{window}(i{+}u, j{+}v, c)\cdot w(u,v,c) + b,
\qquad H_{out} = \left\lfloor \frac{H + 2P - K}{S} \right\rfloor + 1 .
\]
\textbf{Worked example} (SVHN notebook, \S4.2). Take a $5\times5$ patch that contains a vertical stroke (columns
\texttt{0 0 1 1 0}) and a vertical-edge filter with rows $[-1, 0, 1]$: at window $[0,0,1]$ the result is
$(-1\cdot0 + 0\cdot0 + 1\cdot1)\times 3 = 3$ (the stroke's left edge), while window $[1,1,0]$ yields $-3$ (its right
edge). This value is cross-checked against the framework's own convolution output in the notebook.
Setting $K=3$, $S=1$ and \textbf{``same'' padding} ($P=1$) preserves the spatial size, $32\times32 \to 32\times32$.

\subsection{ReLU, pooling, batch normalization}
\begin{table}[H]
\centering
\caption{Layer operations with worked examples (from the notebooks).}
\small
\begin{tabular}{p{3.2cm}p{5.6cm}p{5.8cm}}
\toprule
Layer & Computation & Example \\
\midrule
ReLU & $\max(0, x)$ & $-3 \to 0$, $\ 3 \to 3$, $\ -1.3 \to 0$ \\
Max pooling $2\times2$ & largest value of each $2\times2$ block; halves $H, W$; no parameters & $[[0.1, 0.9],[0.4, 0.2]] \to 0.9$ \\
Batch normalization & $\hat{x} = (x-\mu)/\sqrt{\sigma^2+\epsilon}$, out $= \gamma\hat{x}+\beta$, per channel over the batch & $[1,3,5] \to [-1.22, 0, 1.22]$; dark/bright pixels $[0.05, 0.10, 0.80, 0.90] \to [-1.06, -0.93, 0.87, 1.12]$ \\
Global average pooling & mean of each feature map & $4\times4\times128 \to 128$ \\
Dropout (0.3) & in training, 30\% of values set to 0 & reduces overfitting \\
\bottomrule
\end{tabular}
\end{table}

The \textbf{receptive field} is the region of the input that can affect a single output unit. It expands by
$RF \leftarrow RF + (K-1)\cdot j$ at each layer, with the jump $j$ doubling every time pooling is applied, reaching
22$\times$22 pixels for CNN-3 and 28$\times$28 for CNN-5 (out of a 32$\times$32 image).

\textbf{Parameter counts.} A conv layer contributes $(K\cdot K\cdot C_{in}+1)\cdot F$, e.g.\ the first layer gives
$(3\cdot3\cdot3+1)\cdot32 = 896$; BatchNorm adds $2F$ trainable values ($\gamma,\beta$); a Dense layer adds
$(\text{inputs}+1)\cdot\text{units}$.

\subsection{Output, loss and training}
\begin{itemize}
\item \textbf{Softmax} (multi-class): $p_k = e^{z_k}/\sum_j e^{z_j}$. \textbf{Cross-entropy}: $L = -\ln p_{\text{true}}$.
  For instance, with logits ``3'' $=2$, ``8'' $=1$ and everything else $0$, $p(\text{``3''}) = 0.41$, giving a loss of
  $0.90$ if the true image is a 3 and $1.90$ if it is actually an 8. Before any training the model starts near
  $\ln 10 = 2.30$ (SVHN) or $\ln 43 = 3.76$ (GTSRB).
\item \textbf{Sigmoid} (binary): $p = 1/(1+e^{-z})$. \textbf{Binary cross-entropy}: $L = -[y\ln p + (1-y)\ln(1-p)]$.
  For example, $z=1.2 \Rightarrow p = 0.77$, so the loss is $0.26$ when $y=1$ and $1.47$ when $y=0$.
\item \textbf{Adam}: $w \leftarrow w - \eta\, \hat m/(\sqrt{\hat v}+\epsilon)$, $\eta = 0.001$. An \textbf{epoch} is one
  full sweep through the training batches (256 samples each). \textbf{Early stopping} halts training once 3 epochs pass
  without a new lowest validation loss, and the weights from the best epoch are restored.
\end{itemize}

\subsection{One image through CNN-3, step by step}
Figures~\ref{fig:trace-maps} and~\ref{fig:trace-out} trace a single SVHN test image (a ``3'') as it passes through a
PyTorch CNN-3. These figures come from a CNN-3 trained for only 3 epochs, purely to demonstrate the flow of
information; the actual reported results are taken from the complete runs described in Section~\ref{sec:results}.

\begin{table}[H]
\centering
\caption{What happens to one image in CNN-3 (SVHN).}
\small
\begin{tabular}{p{2.3cm}p{2.4cm}p{9.8cm}}
\toprule
Step & Output shape & What it means \\
\midrule
input & $3\times32\times32$ & 3,072 values in $[0,1]$: the red, green and blue intensity of every pixel \\
block 1 & $32\times16\times16$ & 32 filters each hunt for one small pattern (an edge, a colour change); each resulting
  map marks \emph{where} that pattern turned up, and pooling shrinks the size by half \\
block 2 & $64\times8\times8$ & the filters combine block-1 maps into curves and stroke fragments spanning a wider area \\
block 3 & $128\times4\times4$ & each unit now covers $22\times22$ pixels, enough to capture parts of a digit (a loop, a
  bar) \\
GAP & $128$ & averaging each $4\times4$ map answers ``how strongly is pattern $i$ present anywhere'' \\
Dense 256 $\to$ Dense 10 & $10$ & weighted combinations of the 128 values produce one score (logit) per digit \\
softmax & $10$ & the scores are converted into probabilities that add up to 1 \\
\bottomrule
\end{tabular}
\end{table}

\begin{figure}[H]
\centering
\includegraphics[width=\textwidth,height=0.42\textheight,keepaspectratio]{trace_feature_maps.png}
\caption{Input image alongside the 4 strongest feature maps after each conv block (bright = strong response). The maps
shrink and grow more abstract at each stage: block 1 still resembles the digit, while block 3 only flags \emph{where}
pieces of the digit lie.}
\label{fig:trace-maps}
\end{figure}

\begin{figure}[H]
\centering
\includegraphics[width=\textwidth,height=0.42\textheight,keepaspectratio]{trace_output.png}
\caption{The 128 values produced by global average pooling (left) and the resulting probabilities (right): $p(3) = 0.75$,
$p(5) = 0.21$. The runner-up class ``5'' shares the ``3''s lower curve, and this 5/3 pair also turns up as one of the
most common confusions in the full results.}
\label{fig:trace-out}
\end{figure}

\subsection{Learning: one backpropagation step by hand}
Training nudges every weight slightly so the loss decreases. In PyTorch this is driven by
\code{loss.backward()} (which computes the gradients) followed by \code{optimizer.step()} (which applies the update).
The simplest possible case is a single neuron taking one input:
\[
z = w x + b, \qquad p = \sigma(z) = \frac{1}{1+e^{-z}}, \qquad L = -[\,y\ln p + (1-y)\ln(1-p)\,].
\]
\textbf{Forward pass} with $x = 2$, $w = 0.5$, $b = 0$ and true label $y = 1$:
$z = 1.0$, $p = \sigma(1.0) = 0.731$, $L = -\ln 0.731 = 0.313$.

\textbf{Backward pass} (chain rule). For the sigmoid plus binary cross-entropy pairing the derivative reduces neatly to
$\partial L/\partial z = p - y$, giving
\[
\frac{\partial L}{\partial z} = 0.731 - 1 = -0.269, \qquad
\frac{\partial L}{\partial w} = \frac{\partial L}{\partial z}\cdot\frac{\partial z}{\partial w} = -0.269 \times x = -0.538, \qquad
\frac{\partial L}{\partial b} = -0.269 .
\]
\textbf{Weight update} (gradient descent, learning rate $\eta = 0.1$): $w \leftarrow w - \eta\,\partial L/\partial w$:
\[
w = 0.5 - 0.1\times(-0.538) = 0.554, \qquad b = 0 - 0.1\times(-0.269) = 0.027 .
\]
\textbf{Verification:} plugging in the updated weights gives $z = 0.554\cdot2 + 0.027 = 1.135$, $p = 0.757$,
$L = 0.279$ — the loss dropped from 0.313 to 0.279 and $p$ shifted closer to the true label of 1. Since the gradient
was negative, the weight increased.

\textbf{Inside a CNN} the identical chain rule propagates backward through every layer:
\begin{table}[H]
\centering
\small
\begin{tabular}{p{3cm}p{11.5cm}}
\toprule
Layer & How the gradient passes through \\
\midrule
Softmax + cross-entropy & $\partial L/\partial z_k = p_k - [k = y]$, e.g.\ with $p = (0.75, 0.21, \dots)$ for (``3'', ``5'')
  and true label ``3'': $-0.25$ for ``3'', $+0.21$ for ``5'' \\
Dense & $\partial L/\partial w_{ij} = \partial L/\partial z_j \cdot x_i$ (same rule as the single neuron above) \\
ReLU & the gradient passes through unchanged ($\times 1$) where the input was $> 0$, and is blocked ($\times 0$) where
  it was clipped to 0 \\
Max pooling & the entire gradient flows to whichever input was the maximum; the other three positions receive 0 \\
Convolution & $\partial L/\partial w$ is the sum over every position of (the gradient at that output times the input
  window there): each filter learns from all the locations it was applied to \\
\bottomrule
\end{tabular}
\end{table}
Adam subsequently rescales each weight's update using running averages of its gradients (Section~2.3), and this whole
cycle repeats once per batch across an epoch.

\subsection{One epoch of training}
An \textbf{epoch} covers one full pass over the training set, processed in batches of 256 samples:

\begin{table}[H]
\centering
\small
\begin{tabular}{p{0.6cm}p{4.2cm}p{9.6cm}}
\toprule
Step & Code (PyTorch) & What happens \\
\midrule
1 & \code{model.train()} & switches to training mode: Dropout active, BatchNorm relies on the current batch's $\mu$, $\sigma$ \\
2 & \code{for x, y in train_loader:} & batches are reshuffled into a fresh random order each epoch \\
3 & \code{optimizer.zero_grad()} & wipes out gradients left over from the previous batch \\
4 & \code{loss = criterion(model(x), y)} & runs the forward pass and computes the loss (Sections 2.4, 2.3) \\
5 & \code{loss.backward()} & derives the loss gradient for every weight (Section 2.5) \\
6 & \code{optimizer.step()} & Adam nudges every weight accordingly \\
7 & \code{model.eval()} + validation & measures loss and accuracy on the validation set without altering weights \\
8 & early stopping & saves the weights whenever validation loss hits a new low; stops after 3 epochs with no improvement \\
\bottomrule
\end{tabular}
\end{table}

\textbf{Updates per epoch:} SVHN's 69,502 training images yield 272 weight updates; GTSRB's 36,287 give 142; diabetes'
210,000 rows give 821. \textbf{Early-stopping example} (\run{SV-TF-CNN-3}): validation loss bottomed out at epoch 11,
showed no improvement through epochs 12--14, so training halted after epoch 14 and epoch 11's weights were restored.

\subsection{Training mode vs.\ evaluation mode}
Two layer types act differently depending on whether the model is training or predicting, which is exactly why the
code toggles between \code{model.train()} and \code{model.eval()} (Keras handles this switch automatically inside
\code{fit} and \code{predict}).

\begin{table}[H]
\centering
\small
\begin{tabular}{p{2.6cm}p{5.8cm}p{6.0cm}}
\toprule
Layer & Training & Evaluation (validation, test) \\
\midrule
BatchNorm & normalizes using the \textbf{current batch's} $\mu, \sigma$ and refreshes the running averages:
  $\text{running} = 0.9\cdot\text{running} + 0.1\cdot\text{batch}$
  & relies on the \textbf{stored running} $\mu, \sigma$ instead, so a given image produces the same output no matter
  which batch it lands in \\
Dropout (0.3) & zeroes out a random 30\% of values and scales the rest by $1/0.7$
  (e.g.\ $[1.0, 2.0, 0.5] \to [1.43, 0, 0.71]$) & is inactive: every value passes through unchanged \\
\bottomrule
\end{tabular}
\end{table}
Skipping \code{model.eval()} would leave test-time predictions random (from Dropout) and dependent on batch
composition (from BatchNorm).

\subsection{1-D convolution on a table}
Each patient is represented as a vector of $d$ standardized features, treated as a length-$d$ sequence with a single
channel. A Conv1D filter of size 3 blends three adjacent features: $\text{out}[i] = w_1x[i-1] + w_2x[i] + w_3x[i+1] + b$.
For example, $x = [0.5, -1.2, 0.3, 2.0, -0.4, 0.0]$ with $w = [1, 0, -1]$ gives
out $= [1.2, 0.2, -3.2, 0.7, 2.0, -0.4]$ (verified against the framework's own Conv1D). Table columns, unlike pixels,
carry no inherent ordering, so which features count as ``neighbours'' depends on an arbitrary column arrangement; for
this reason the tabular models \textbf{flatten} the final feature map rather than averaging it (so that position
still identifies which feature it is), and an MLP baseline is trained for comparison.

\subsection{Metrics}
For each class: precision $= TP/(TP+FP)$, recall $= TP/(TP+FN)$, $F1 = 2PR/(P+R)$. \textbf{Macro F1} takes the
unweighted average of per-class F1 scores. Accuracy can be deceptive when classes are imbalanced: in a toy test set of
95 ``big class'' images and 5 ``rare class'' images, always guessing the big class scores accuracy 0.95 but only macro
F1 0.49. \textbf{ROC-AUC} (for binary problems) measures the chance that a randomly chosen positive sample receives a
higher $p$ than a randomly chosen negative one (0.5 corresponds to random guessing).

% =====================================================================
\section{Project Structure and Setup}
% =====================================================================

\begin{table}[H]
\centering
\caption{Folders and naming convention.}
\small
\begin{tabular}{p{4.4cm}p{10.4cm}}
\toprule
Item & Content / rule \\
\midrule
\code{dataset/sv, gt, db} & raw Kaggle files (not committed) + \code{splits.npz} (row indices of the split) \\
\code{notebook/} & \code{sv_svhn.ipynb}, \code{gt_gtsrb.ipynb}, \code{db_diabetes.ipynb}: one self-contained notebook per dataset \\
\code{report/} & this report, \code{images/} (figures extracted from the notebooks), \code{make_summary.py} \\
Dataset / framework / model & \run{SV}, \run{GT}, \run{DB} / \run{TF}, \run{PT} / \run{CNN-3}, \run{CNN-5} (number of conv layers) \\
Run ID and figure name & \run{<dataset>-<framework>-<model>}; figure \code{<run_id>_<plot>.png} \\
\bottomrule
\end{tabular}
\end{table}

Every notebook follows the same layout: setup, data, exploration, TensorFlow, PyTorch, comparison, conclusion.
Each code cell is preceded by a brief explanatory cell covering the formula, a worked example, or the relevant shapes.

\textbf{Rules for a fair comparison.} A stratified 70/15/15 split is computed once and reused by both frameworks;
batch size 256; Adam ($\eta = 10^{-3}$, $\epsilon = 10^{-7}$); a cap of 20 epochs; early-stopping patience of 3;
Glorot-uniform weight initialization with zero biases in both frameworks (PyTorch is re-initialized to match Keras);
BatchNorm momentum of 0.9 in Keras corresponds to 0.1 in PyTorch (both express
$\text{running} = 0.9\cdot\text{running} + 0.1\cdot\text{batch}$), with $\epsilon = 10^{-5}$; each PyTorch model's
trainable parameter count is asserted to match its Keras counterpart exactly.

\textbf{Environment for the reported run.} Windows, Python 3.13, TensorFlow 2.21, PyTorch 2.14 (CPU build).
Both frameworks were run on the \textbf{CPU} (native Windows TensorFlow has no GPU support, and the PyTorch build used
was CPU-only), so all speed comparisons are CPU-to-CPU. Per-epoch timing is the median across epoch 2 onward, since
epoch 1 includes setup overhead.

% =====================================================================
\section{Datasets}
% =====================================================================

\begin{table}[H]
\centering
\caption{Datasets (all from Kaggle, download $<$ 300 MB each).}
\small
\begin{tabular}{p{2.0cm}p{5.0cm}p{2.2cm}p{1.2cm}p{3.6cm}}
\toprule
Code & Source & Samples & Classes & Input \\
\midrule
\run{SV} SVHN & {\footnotesize\url{sahityasetu/street-view-house-numbers-images}} & 99,289 & 10 & $32\times32\times3$ photo \\
\run{GT} GTSRB & {\footnotesize\url{harbhajansingh21/german-traffic-sign-dataset}} & 51,839 & 43 & $32\times32\times3$ photo \\
\run{DB} Diabetes & Playground Series S5E12 (\code{train.csv}) & 300,000 of 700,000 & 2 & 24 features $\to d = 42$ \\
\bottomrule
\end{tabular}
\end{table}

\subsection{SVHN}
Cropped digit images taken from Google Street View, where the label refers only to the centre digit and any
neighbouring digits are treated as noise. The two MATLAB files are read via \code{scipy.io.loadmat}, reshaped from
$(32,32,3,N)$ to $(N,32,32,3)$, and label 10 is remapped to digit 0. The classes are imbalanced: digit 1 accounts for
13,272 images versus 4,378 for digit 9 (a 3.0 ratio). Pixel values are rescaled via $x/255$ ($0\to0$, $128\to0.50$,
$255\to1$); Keras consumes $(N,32,32,3)$ while PyTorch consumes $(N,3,32,32)$.

\begin{figure}[H]
\centering
\includegraphics[width=\textwidth,height=0.42\textheight,keepaspectratio]{SV_samples.png}
\par\medskip
\includegraphics[width=\textwidth,height=0.42\textheight,keepaspectratio]{SV_class_distribution.png}
\caption{\run{SV}: 8 training images per digit (top) and images per class and split (bottom).}
\end{figure}

\subsection{GTSRB}
Traffic-sign photographs shot from a moving vehicle and resized to $32\times32$. The three pickle files supplied are
merged and then re-split. Class imbalance is pronounced (class 0, ``Speed limit 20'', has just 189 images versus 2,100
for class 2, ``Speed limit 50'', an 11.1 ratio), which is why macro F1 serves as the headline metric. A large share of
the photos are quite dark, an effect that BatchNorm helps offset.

\begin{figure}[H]
\centering
\includegraphics[width=\textwidth,height=0.42\textheight,keepaspectratio]{GT_samples.png}
\par\medskip
\includegraphics[width=\textwidth,height=0.42\textheight,keepaspectratio]{GT_class_distribution.png}
\caption{\run{GT}: one training image per class (top) and images per class and split (bottom).}
\end{figure}

\subsection{Diabetes (Playground Series S5E12)}
This dataset holds 700,000 synthetic patient records generated by Kaggle from a real-world source, from which a
stratified sample of 300,000 is drawn. Every row is converted into $d = 42$ numbers: 15 numeric features standardized
using the training statistics ($z = (x-\mu)/\sigma$, e.g.\ BMI 31.7 $\to (31.7-25.9)/2.9 = 2.0$), 6 categorical
features one-hot encoded, and 3 binary flags retained as-is. Since 62.3\% of patients are diabetic, class weights
$w_c = N/(2N_c)$ are applied: $w_0 = 1.327$, $w_1 = 0.802$. The most discriminating numeric features (age, physical
activity) separate the two classes by only about 0.4 standard deviations, an early hint that this will be a difficult
prediction task.

\begin{figure}[H]
\centering
\includegraphics[width=0.7\textwidth,height=0.42\textheight,keepaspectratio]{DB_feature_effect.png}
\par\medskip
\includegraphics[width=\textwidth,height=0.42\textheight,keepaspectratio]{DB_top_features.png}
\caption{\run{DB}: standardized mean difference per numeric feature (top) and the two most separating features (bottom).}
\end{figure}

% =====================================================================
\section{Models}
% =====================================================================

\begin{table}[H]
\centering
\caption{Image models (SVHN, GTSRB). 3$\times$3 filters, padding ``same''; each Conv is followed by BatchNorm and ReLU.}
\footnotesize
\setlength{\tabcolsep}{4pt}\begin{tabular}{lll}
\toprule
Step & CNN-3 output & CNN-5 output \\
\midrule
input & $32\times32\times3$ & $32\times32\times3$ \\
block 1 & Conv 32 $\to$ Pool $\to 16\times16\times32$ & Conv 32 $\to$ Conv 32 $\to$ Pool $\to 16\times16\times32$ \\
block 2 & Conv 64 $\to$ Pool $\to 8\times8\times64$ & Conv 64 $\to$ Conv 64 $\to$ Pool $\to 8\times8\times64$ \\
block 3 & Conv 128 $\to$ Pool $\to 4\times4\times128$ & Conv 128 $\to$ Pool $\to 4\times4\times128$ \\
head & GAP $\to$ Dropout 0.3 $\to$ Dense 256 $\to$ Dense 10/43 & same \\
receptive field & 22$\times$22 & 28$\times$28 \\
\bottomrule
\end{tabular}
\end{table}

\begin{table}[H]
\centering
\caption{Tabular models (diabetes). Conv1D kernel 3, padding ``same'', no pooling.}
\footnotesize
\setlength{\tabcolsep}{4pt}\begin{tabular}{lll}
\toprule
Step & CNN-3 & CNN-5 \\
\midrule
Conv1D filters & 32 $\to$ 64 $\to$ 64 ($d\times64$) & 32 $\to$ 32 $\to$ 64 $\to$ 64 $\to$ 128 ($d\times128$) \\
head & Flatten $\to$ Dropout 0.3 $\to$ Dense 64 $\to$ Dense 1 & same \\
features one position sees & 7 & 11 \\
\bottomrule
\end{tabular}
\end{table}

\begin{table}[H]
\centering
\caption{Trainable parameters (identical in TF and PT).}
\small
\begin{tabular}{lrrr}
\toprule
Model & SVHN & GTSRB & Diabetes \\
\midrule
CNN-3 & 129,290 & 137,771 & 191,169 \\
CNN-5 & 175,658 & 184,139 & 391,329 \\
MLP baseline & -- & -- & 6,977 \\
\bottomrule
\end{tabular}
\end{table}

\subsection{TensorFlow, Keras and PyTorch}
\begin{itemize}
\item \textbf{TensorFlow} (Google) is the lower-level library: it handles tensors, automatic differentiation, and
  CPU/GPU execution.
\item \textbf{Keras} sits \emph{inside} TensorFlow as its high-level API — a model is declared layer by layer and
  trained with \code{compile} + \code{fit}, while TensorFlow does the actual computing underneath.
\item \textbf{PyTorch} (Meta) is a separate library that operates at roughly the same level as TensorFlow itself.
\end{itemize}
In other words, ``TensorFlow'' and ``Keras'' name one side of this comparison (models written in Keras, run by
TensorFlow, run ID \run{TF}), and PyTorch names the other (run ID \run{PT}).

\begin{table}[H]
\centering
\caption{How the same CNN is written and trained in the two frameworks.}
\small
\begin{tabular}{p{2.6cm}p{5.8cm}p{6.2cm}}
\toprule
 & Keras / TensorFlow & PyTorch \\
\midrule
model & layers stacked declaratively: \code{Conv2D(32, 3, padding="same")} & \code{nn.Module} subclass: \code{nn.Conv2d(3, 32, 3, padding=1)}, input channels specified manually \\
image axes & $(N, 32, 32, 3)$, channels last & $(N, 3, 32, 32)$, channels first \\
training & single \code{compile(...)} + \code{fit(...)} call & a manually written loop: \code{zero_grad} $\to$ forward $\to$ loss $\to$ \code{backward} $\to$ \code{step} \\
early stopping & handled by the built-in \code{EarlyStopping} callback & tracked with a counter inside the loop \\
train / test mode & toggled automatically by \code{fit} / \code{predict} & \code{model.train()} / \code{model.eval()} invoked explicitly \\
device & picks up a GPU automatically when present & requires \code{.to(device)} on the model and every batch \\
strength & concise, fast to write & every step is explicit and easy to modify \\
\midrule
results here & \multicolumn{2}{p{12cm}}{comparable accuracy (gap under 0.011 with matching architecture and settings); in this CPU run Keras finished 9--16\% faster per epoch on the image datasets and about 50\% faster on the tabular one} \\
\bottomrule
\end{tabular}
\end{table}

\subsection{Implementation (TensorFlow/Keras)}
\begin{table}[H]
\centering
\caption{Keras building blocks used in the image models.}
\small
\begin{tabular}{p{3.0cm}p{11.2cm}}
\toprule
Step & Keras / TensorFlow \\
\midrule
data & \code{tf.data.Dataset}: shuffle $\to$ \code{batch(256)} $\to$ divide by 255 $\to$ \code{prefetch}; images $(N, 32, 32, 3)$ \\
conv block & \code{Conv2D(F, 3, padding="same")} $\to$\newline \code{BatchNormalization(momentum=0.9)} $\to$ \code{ReLU()} $\to$ \code{MaxPooling2D(2)} \\
head & \code{GlobalAveragePooling2D()} $\to$ \code{Dropout(0.3)} $\to$ \code{Dense(256, activation="relu")} $\to$ \code{Dense(n_classes)} \\
init & Glorot uniform weights, zero biases (Keras default) \\
loss / optimizer & \code{SparseCategoricalCrossentropy(from_logits=True)} (diabetes: \code{BinaryCrossentropy}), \code{Adam(1e-3)} \\
early stopping & \code{EarlyStopping(monitor="val_loss", patience=3,}\newline \code{restore_best_weights=True)} \\
\bottomrule
\end{tabular}
\end{table}

\noindent\begin{minipage}{\linewidth}
\begin{lstlisting}
def build_tf(arch, input_shape, n_classes):
    inputs = keras.Input(shape=input_shape)             # (32, 32, 3)
    x = inputs
    for filters, pool in arch:                          # e.g. [(32, True), (64, True), (128, True)]
        x = layers.Conv2D(filters, 3, padding="same")(x)
        x = layers.BatchNormalization(momentum=0.9, epsilon=1e-5)(x)
        x = layers.ReLU()(x)
        if pool:
            x = layers.MaxPooling2D(2)(x)
    x = layers.GlobalAveragePooling2D()(x)
    x = layers.Dropout(0.3)(x)
    x = layers.Dense(256, activation="relu")(x)
    return keras.Model(inputs, layers.Dense(n_classes)(x))    # logits

model.compile(optimizer=keras.optimizers.Adam(1e-3),
              loss=keras.losses.SparseCategoricalCrossentropy(from_logits=True),
              metrics=["accuracy"])
model.fit(train_ds, validation_data=val_ds, epochs=20,
          callbacks=[keras.callbacks.EarlyStopping(monitor="val_loss", patience=3,
                                                   restore_best_weights=True)])
\end{lstlisting}
\end{minipage}

A single call to \code{fit} executes the entire training loop from Section 2.6 (batching, forward pass, backward pass,
Adam update, validation, early stopping); the PyTorch version spells these same steps out explicitly (see below).

\subsection{Implementation (PyTorch)}
Here is the equivalent model built in PyTorch. Its layers and settings match the Keras version exactly (identical
parameter counts, confirmed in the notebooks), but the training loop itself is written by hand.

\begin{table}[H]
\centering
\caption{PyTorch building blocks used in the image models.}
\small
\begin{tabular}{p{3.0cm}p{11.2cm}}
\toprule
Step & PyTorch \\
\midrule
data & \code{TensorDataset} + \code{DataLoader(batch_size=256, shuffle=True)}; images $(N, 3, 32, 32)$, divided by 255 per batch \\
conv block & \code{Conv2d(C_in, F, 3, padding=1)} $\to$ \code{BatchNorm2d(F)} $\to$ \code{ReLU()} $\to$ \code{MaxPool2d(2)} \\
head & \code{AdaptiveAvgPool2d(1)} $\to$ \code{Flatten()} $\to$ \code{Dropout(0.3)} $\to$ \code{Linear(128, 256)} $\to$ \code{ReLU()} $\to$ \code{Linear(256, n_classes)} \\
init & \code{xavier_uniform_} weights, zero biases \\
loss / optimizer & \code{CrossEntropyLoss()} (diabetes: \code{BCEWithLogitsLoss}), \code{Adam(lr=1e-3)} \\
early stopping & keep the best \code{state_dict}; stop after 3 epochs without a lower validation loss \\
\bottomrule
\end{tabular}
\end{table}

\noindent\begin{minipage}{\linewidth}
\begin{lstlisting}
class CNN(nn.Module):
    def __init__(self, arch, in_channels, n_classes):
        super().__init__()
        blocks, c = [], in_channels
        for filters, pool in arch:                 # e.g. [(32, True), (64, True), (128, True)]
            blocks += [nn.Conv2d(c, filters, 3, padding=1), nn.BatchNorm2d(filters), nn.ReLU()]
            if pool:
                blocks.append(nn.MaxPool2d(2))
            c = filters
        self.features = nn.Sequential(*blocks)
        self.head = nn.Sequential(nn.AdaptiveAvgPool2d(1), nn.Flatten(), nn.Dropout(0.3),
                                  nn.Linear(c, 256), nn.ReLU(), nn.Linear(256, n_classes))

    def forward(self, x):
        return self.head(self.features(x))

for x, y in train_loader:                          # one epoch
    optimizer.zero_grad()                          # clear gradients
    loss = criterion(model(x), y)                  # forward + cross-entropy
    loss.backward()                                # backpropagation
    optimizer.step()                               # Adam update
\end{lstlisting}
\end{minipage}

% =====================================================================
\section{Results per Dataset}\label{sec:results}
% =====================================================================

\subsection*{How to read the figures}
\begin{itemize}
\item \textbf{Training curves} plot loss and accuracy per epoch for both the training and validation sets. Both curves
  falling together means the model is still learning; training loss falling while validation loss climbs is
  \textbf{overfitting}, and early stopping stops the run at the epoch just before that starts.
\item \textbf{Confusion matrix} (row-normalized). Rows are the true class, columns the predicted one, and each row
  sums to 1, so the diagonal entries are each class's recall. In Figure~\ref{fig:sv-conf}, for instance, row 7 /
  column 1 = 0.032 means about 3\% of the real 7s got classified as 1.
\item \textbf{Feature maps} are grayscale images of one filter's output — bright pixels show where that filter's
  pattern was detected.
\item \textbf{ROC curve} (diabetes) plots true-positive rate against false-positive rate across all thresholds. The
  diagonal is random guessing (AUC 0.5); the top-left corner is a perfect classifier (AUC 1).
\end{itemize}

\subsection{SVHN}
\begin{table}[H]
\centering
\caption{\run{SV} test results.}
\small
\begin{tabular}{lrrrrrr}
\toprule
Run ID & Accuracy & Macro F1 & Epochs & s/epoch & Train time (s) & ms / 1k pred. \\
\midrule
\run{SV-TF-CNN-3} & 0.9012 & 0.8959 & 14 & 19.8 & 279 & 0.07 \\
\run{SV-TF-CNN-5} & \textbf{0.9398} & \textbf{0.9363} & 20 & 42.9 & 865 & 0.14 \\
\run{SV-PT-CNN-3} & 0.9091 & 0.9069 & 20 & 28.1 & 557 & 0.15 \\
\run{SV-PT-CNN-5} & 0.9319 & 0.9282 & 12 & 61.5 & 750 & 0.27 \\
\bottomrule
\end{tabular}
\end{table}

\begin{figure}[H]
\centering
\includegraphics[width=\textwidth,height=0.42\textheight,keepaspectratio]{SV_comparison_curves.png}
\caption{\run{SV}: validation loss and accuracy of the four runs (solid = TF, dashed = PT).}
\end{figure}

\begin{figure}[H]
\centering
\includegraphics[width=0.62\textwidth,height=0.42\textheight,keepaspectratio]{SV-TF-CNN-5_confusion.png}
\par\medskip
\includegraphics[width=\textwidth,height=0.42\textheight,keepaspectratio]{SV-TF-CNN-5_misclassified.png}
\caption{\run{SV-TF-CNN-5}: confusion matrix (top) and misclassified test images (bottom).}\label{fig:sv-conf}
\end{figure}

\begin{figure}[H]
\centering
\includegraphics[width=\textwidth,height=0.42\textheight,keepaspectratio]{SV-TF-CNN-3_feature_maps.png}
\par\medskip
\includegraphics[width=\textwidth,height=0.42\textheight,keepaspectratio]{SV-TF-CNN-5_feature_maps.png}
\caption{First-layer feature maps after ReLU for one test image: \run{SV-TF-CNN-3} (top), \run{SV-TF-CNN-5} (bottom).}
\end{figure}

\begin{itemize}
\item \textbf{Depth.} Moving to CNN-5 raises macro F1 by +0.040 (TF) and +0.021 (PT), and accuracy by +0.039 (TF) and
  +0.023 (PT). Under PyTorch it also reaches its best validation loss sooner (epoch 9 versus 19), though every epoch
  costs roughly twice as much time either way.
\item \textbf{Framework.} At CNN-3, TF trails PT by about 0.008 accuracy; at CNN-5 that flips and TF leads PT by about
  the same margin.
\item \textbf{Errors.} On the best model (\run{SV-TF-CNN-5}), the most common mistake is a 4 read as a 1 (4.6\% of the
  4s), then 5 read as 3 (3.4\%) and 7 read as 1 (3.2\%) — all digits that share a vertical stroke. Looking through the
  misclassified-image grid, the usual suspects show up: blurry crops, low contrast, and digits crowded by their
  neighbours.
\item The first-layer feature maps mostly respond to the digit's strokes and edges, though a few pick up on
  background colour instead.
\end{itemize}

\subsection{GTSRB}
\begin{table}[H]
\centering
\caption{\run{GT} test results.}
\small
\begin{tabular}{lrrrrrr}
\toprule
Run ID & Accuracy & Macro F1 & Epochs & s/epoch & Train time (s) & ms / 1k pred. \\
\midrule
\run{GT-TF-CNN-3} & 0.9801 & 0.9745 & 20 & 10.2 & 207 & 0.07 \\
\run{GT-TF-CNN-5} & 0.9924 & 0.9893 & 12 & 21.9 & 268 & 0.15 \\
\run{GT-PT-CNN-3} & 0.9685 & 0.9621 & 14 & 14.2 & 199 & 0.15 \\
\run{GT-PT-CNN-5} & \textbf{0.9973} & \textbf{0.9979} & 19 & 30.0 & 572 & 0.26 \\
\bottomrule
\end{tabular}
\end{table}

\begin{figure}[H]
\centering
\includegraphics[width=\textwidth,height=0.42\textheight,keepaspectratio]{GT_comparison_curves.png}
\caption{\run{GT}: validation loss and accuracy of the four runs (solid = TF, dashed = PT).}
\end{figure}

\begin{figure}[H]
\centering
\includegraphics[width=0.85\textwidth,height=0.42\textheight,keepaspectratio]{GT-PT-CNN-5_confusion.png}
\par\medskip
\includegraphics[width=\textwidth,height=0.42\textheight,keepaspectratio]{GT-PT-CNN-5_misclassified.png}
\caption{\run{GT-PT-CNN-5}: confusion matrix over 43 classes (top) and misclassified test signs (bottom).}
\end{figure}

\begin{figure}[H]
\centering
\includegraphics[width=0.75\textwidth,height=0.42\textheight,keepaspectratio]{GT-TF-CNN-5_filters.png}
\par\medskip
\includegraphics[width=\textwidth,height=0.42\textheight,keepaspectratio]{GT-TF-CNN-5_feature_maps.png}
\caption{\run{GT-TF-CNN-5}: the 32 first-layer filters (top) and feature maps for one test sign (bottom).}
\end{figure}

\begin{itemize}
\item \textbf{Depth.} CNN-5 raises macro F1 by +0.015 (TF) and +0.036 (PT). Macro F1 weighs all 43 classes equally, so
  this gain is not just coming from the common speed-limit signs — the rare sign types improve too.
\item \textbf{Framework.} The winner flips with depth. At CNN-3, TF leads PT by 0.012 macro F1 (TF used the full 20
  epochs to get there; PT stopped at epoch 14). At CNN-5, PT leads TF by 0.009 macro F1, and \run{GT-PT-CNN-5} ends up
  the strongest model in the whole study.
\item The confusion matrix is almost entirely diagonal. What the misclassified-image grid shows is mostly dark or
  blurred photos — the remaining hard cases.
\end{itemize}

\subsection{Diabetes}
\begin{table}[H]
\centering
\caption{\run{DB} test results (threshold 0.5).}
\small
\begin{tabular}{lrrrrrrr}
\toprule
Run ID & Accuracy & Macro F1 & F1 (pos.) & ROC-AUC & Params & Epochs & s/epoch \\
\midrule
\run{DB-TF-CNN-3} & 0.6452 & \textbf{0.6304} & 0.7045 & \textbf{0.6963} & 191,169 & 7 & 9.7 \\
\run{DB-TF-CNN-5} & 0.6358 & 0.6275 & 0.6831 & 0.6945 & 391,329 & 5 & 18.6 \\
\run{DB-PT-CNN-3} & 0.6400 & 0.6288 & 0.6934 & 0.6951 & 191,169 & 8 & 21.1 \\
\run{DB-PT-CNN-5} & \textbf{0.6454} & 0.6303 & \textbf{0.7051} & 0.6955 & 391,329 & 6 & 38.4 \\
\run{DB-TF-MLP} (baseline) & 0.6276 & 0.6221 & 0.6677 & 0.6946 & 6,977 & 5 & 1.2 \\
\bottomrule
\end{tabular}
\end{table}

\begin{figure}[H]
\centering
\includegraphics[width=\textwidth,height=0.42\textheight,keepaspectratio]{DB_comparison_curves.png}
\caption{\run{DB}: validation loss and accuracy of the four runs (solid = TF, dashed = PT).}
\end{figure}

\begin{figure}[H]
\centering
\includegraphics[width=\textwidth,height=0.42\textheight,keepaspectratio]{DB-TF_roc_pr.png}
\par\medskip
\includegraphics[width=0.65\textwidth,height=0.42\textheight,keepaspectratio]{DB-TF-CNN-5_probabilities.png}
\caption{\run{DB-TF}: ROC and precision-recall curves (top); predicted probability by true class for \run{DB-TF-CNN-5} (bottom).}
\end{figure}

\begin{itemize}
\item \textbf{Depth.} There is no consistent benefit here. Going from CNN-3 to CNN-5, ROC-AUC moves by $-0.0018$ (TF)
  and $+0.0004$ (PT), despite the parameter count doubling and per-epoch time growing by 1.8--1.9$\times$.
\item \textbf{Framework.} TF and PT stay within 0.0012 ROC-AUC of one another. The accuracy differences that do show up
  ($\pm0.01$) just reflect where individual predictions happen to fall around the 0.5 threshold.
\item \textbf{Baseline.} An MLP with 56$\times$ fewer parameters still reaches ROC-AUC 0.6946, about the same as the
  CNNs. The convolution's core assumption — that neighbouring features carry information — does not seem to hold
  here.
\item The validation curves oscillate within a narrow band (loss 0.62--0.645). With such a weak signal, every model
  settles into a noisy optimum after only 2--5 epochs, and early stopping cuts training short soon after that.
\item The predicted-probability histogram shows the two classes overlapping heavily. Moving the decision threshold
  away from 0.5 would just trade precision for recall, not actually separate the classes.
\end{itemize}

% =====================================================================
\section{Cross-Dataset Comparison}
% =====================================================================

\begin{figure}[H]
\centering
\includegraphics[width=\textwidth,height=0.42\textheight,keepaspectratio]{summary_f1.png}
\par\medskip
\includegraphics[width=\textwidth,height=0.42\textheight,keepaspectratio]{summary_time.png}
\caption{All 12 runs: test macro F1 (top) and training time per epoch (bottom). Purple = CNN-3, green = CNN-5; the hatched bars are PyTorch, the plain bars are TensorFlow/Keras.}
\end{figure}

\begin{table}[H]
\centering
\caption{Effect of depth (CNN-5 minus CNN-3) and of framework (PT minus TF).}
\small
\begin{tabular}{llrrr}
\toprule
Comparison & & SVHN & GTSRB & Diabetes \\
\midrule
CNN-5 $-$ CNN-3, macro F1 & TF & +0.0404 & +0.0148 & $-0.0029$ \\
 & PT & +0.0213 & +0.0358 & +0.0015 \\
CNN-5 / CNN-3, time per epoch & TF & 2.17$\times$ & 2.15$\times$ & 1.92$\times$ \\
 & PT & 2.19$\times$ & 2.11$\times$ & 1.82$\times$ \\
\midrule
PT $-$ TF, macro F1 & CNN-3 & +0.0110 & $-0.0124$ & $-0.0016$ \\
 & CNN-5 & $-0.0081$ & +0.0086 & +0.0028 \\
PT / TF, time per epoch & CNN-3 & 1.42$\times$ & 1.39$\times$ & 2.18$\times$ \\
 & CNN-5 & 1.43$\times$ & 1.37$\times$ & 2.06$\times$ \\
\bottomrule
\end{tabular}
\end{table}

\begin{itemize}
\item Depth pays off wherever the input has spatial structure — +0.01 to +0.04 macro F1 on both image datasets — but
  changes nothing on the tabular data, where the ceiling is set by the features themselves and Conv1D has no real
  spatial structure to exploit.
\item Both frameworks reach comparable quality: every PT--TF gap stays under 0.013 macro F1, about the same size as
  the noise you'd expect from run to run with a single seed — though which framework comes out ahead is not
  consistent, flipping between CNN-3 and CNN-5 on both image datasets. Where they clearly differ is speed. On this
  CPU run, Keras was quicker per epoch across the board, and the gap was largest on the tabular model, where
  PyTorch's per-batch Python loop adds the most overhead (roughly 2$\times$ slower there versus 1.4$\times$ on the
  image datasets). Keras also predicted faster at inference time — roughly 1.7--2.1$\times$ per 1,000 images on the
  image datasets, 1.3--1.5$\times$ on the tabular data.
\item CNN-5 costs roughly twice the time per epoch, and because it does not consistently converge in fewer epochs on
  this run, its total training time ends up anywhere from 1.3$\times$ to 3.1$\times$ that of CNN-3 on the image
  datasets, depending on the run.
\end{itemize}

\subsection{Keras vs.\ PyTorch outputs}
The same model trained under both frameworks, placed side by side (test set).

\begin{table}[H]
\centering
\small
\begin{tabular}{llrrrrrrrr}
\toprule
 & & \multicolumn{4}{c}{Keras / TensorFlow} & \multicolumn{4}{c}{PyTorch} \\
\cmidrule(lr){3-6}\cmidrule(lr){7-10}
Dataset & Model & Acc. & Macro F1 & Epochs & s/epoch & Acc. & Macro F1 & Epochs & s/epoch \\
\midrule
SVHN & CNN-3 & 0.9012 & 0.8959 & 14 & 19.8 & 0.9091 & 0.9069 & 20 & 28.1 \\
SVHN & CNN-5 & 0.9398 & 0.9363 & 20 & 42.9 & 0.9319 & 0.9282 & 12 & 61.5 \\
GTSRB & CNN-3 & 0.9801 & 0.9745 & 20 & 10.2 & 0.9685 & 0.9621 & 14 & 14.2 \\
GTSRB & CNN-5 & 0.9924 & 0.9893 & 12 & 21.9 & 0.9973 & 0.9979 & 19 & 30.0 \\
Diabetes & CNN-3 & 0.6452 & 0.6304 & 7 & 9.7 & 0.6400 & 0.6288 & 8 & 21.1 \\
Diabetes & CNN-5 & 0.6358 & 0.6275 & 5 & 18.6 & 0.6454 & 0.6303 & 6 & 38.4 \\
\bottomrule
\end{tabular}
\end{table}

\begin{figure}[H]
\centering
\includegraphics[width=\textwidth,height=0.40\textheight,keepaspectratio]{SV-TF-CNN-5_curves.png}
\par\medskip
\includegraphics[width=\textwidth,height=0.40\textheight,keepaspectratio]{SV-PT-CNN-5_curves.png}
\caption{Training curves of CNN-5 on SVHN: \run{SV-TF-CNN-5} (top) and \run{SV-PT-CNN-5} (bottom).}
\end{figure}
\begin{figure}[H]
\centering
\includegraphics[width=0.62\textwidth,height=0.40\textheight,keepaspectratio]{SV-TF-CNN-5_confusion.png}
\par\medskip
\includegraphics[width=0.62\textwidth,height=0.40\textheight,keepaspectratio]{SV-PT-CNN-5_confusion.png}
\caption{Confusion matrices of CNN-5 on SVHN: \run{SV-TF-CNN-5} (top) and \run{SV-PT-CNN-5} (bottom).}
\end{figure}
\begin{figure}[H]
\centering
\includegraphics[width=\textwidth,height=0.40\textheight,keepaspectratio]{GT-TF-CNN-5_curves.png}
\par\medskip
\includegraphics[width=\textwidth,height=0.40\textheight,keepaspectratio]{GT-PT-CNN-5_curves.png}
\caption{Training curves of CNN-5 on GTSRB: \run{GT-TF-CNN-5} (top) and \run{GT-PT-CNN-5} (bottom).}
\end{figure}
\begin{figure}[H]
\centering
\includegraphics[width=\textwidth,height=0.40\textheight,keepaspectratio]{DB-TF_roc_pr.png}
\par\medskip
\includegraphics[width=\textwidth,height=0.40\textheight,keepaspectratio]{DB-PT_roc_pr.png}
\caption{ROC and precision-recall curves on diabetes: Keras models (top) and PyTorch models (bottom).}
\end{figure}

\subsection{Why depth helps on images: the receptive field}
Every $3\times3$ convolution extends a unit's view by one pixel on each side, and every pooling operation doubles the
stride between neighbouring units. Consequently the \textbf{receptive field} (the input region visible to one unit)
grows as $RF \leftarrow RF + (K-1)\cdot j$, with the jump $j$ doubling each time pooling occurs:

\begin{table}[H]
\centering
\small
\begin{tabular}{lcccccccc}
\toprule
CNN-3 & conv1 & +pool & conv2 & +pool & conv3 & +pool & & \\
receptive field & 3 & 4 & 8 & 10 & 18 & \textbf{22} & & \\
\midrule
CNN-5 & conv1 & conv2 & +pool & conv3 & conv4 & +pool & conv5 & +pool \\
receptive field & 3 & 5 & 6 & 10 & 14 & 16 & 24 & \textbf{28} \\
\bottomrule
\end{tabular}
\end{table}

\begin{figure}[H]
\centering
\includegraphics[width=\textwidth,height=0.42\textheight,keepaspectratio]{receptive_field.png}
\caption{Receptive field after each block, drawn on an SVHN image. A final CNN-3 unit sees $22\times22$ pixels; a
CNN-5 unit sees $28\times28$, almost the whole digit and its neighbours.}
\end{figure}

Because of this, a CNN-5 unit can connect the top and bottom of a digit (for example, the open or closed left side
distinguishing a 3 from a 5 or an 8), and it also benefits from two additional non-linear layers for combining
patterns. This lines up with the observed results: +0.02 to +0.04 macro F1 on SVHN, where digits occupy most of the
frame and have visual look-alikes, and +0.015 to +0.036 on GTSRB.

\subsection{Why a CNN does not help on the diabetes table}
\begin{table}[H]
\centering
\small
\begin{tabular}{p{3.6cm}p{5.2cm}p{5.6cm}}
\toprule
 & Image (SVHN, GTSRB) & Table (diabetes) \\
\midrule
neighbours & adjacent pixels typically belong to the same stroke or edge & adjacent columns are unrelated to one
  another (e.g.\ \code{screen_time} sitting next to \code{bmi}); the ordering was simply our choice \\
same pattern, other place & an edge looks the same regardless of location, so one filter can be reused across the
  whole image & here position \emph{is} the feature; a pattern found at ``age'' means something entirely different at
  ``cholesterol'' \\
signal in the data & strong: models reach 0.90--1.00 & weak: the top feature separates the classes by only
  $\approx$0.4 standard deviations \\
result & CNN-5 $>$ CNN-3 & CNN-3 $\approx$ CNN-5 $\approx$ MLP (ROC-AUC 0.694--0.696) \\
\bottomrule
\end{tabular}
\end{table}
The two premises that make convolution effective for images — nearby pixels matter, and the same filter is reusable
everywhere — simply do not hold for tabular data. Adding depth or more filters therefore just adds parameters without
adding capability; the true ceiling is set by how much information the 24 features actually contain, and every model
arrives at that ceiling within 2--5 epochs.

% =====================================================================
\section{Conclusion and Limitations}
% =====================================================================

\begin{enumerate}
\item Going from CNN-3 to CNN-5 improves SVHN and GTSRB by 1.5 to 4 macro-F1 points. The extra two layers widen the
  receptive field from 22$\times$22 to 28$\times$28, at roughly double the cost per epoch.
\item On the diabetes table, neither extra depth nor convolution itself helps — all four CNNs and even a
  7k-parameter MLP land around ROC-AUC 0.695. A CNN is not a natural fit for this kind of tabular data.
\item The same architecture gives a broadly similar result in TensorFlow/Keras and PyTorch, with macro-F1 gaps under
  0.013 on every dataset, though which framework comes out ahead is not consistent. That is compatible with the two
  frameworks differing mainly in how much they abstract away rather than in the underlying learning process, though a
  single seed per run cannot fully separate a real framework effect from ordinary training noise.
\item Keras needed less code to write and ran faster in this CPU setup; PyTorch trades that away for a training loop
  where every step is visible and easy to modify.
\end{enumerate}

\textbf{Limitations.} Each run used a single seed (so no variance estimate is available); GPU/CPU kernels are not
fully deterministic; no data augmentation was applied; the diabetes dataset is synthetic; Conv1D's results depend on
an arbitrary column ordering; all timings come from one CPU-only Windows machine; and three of the twelve image-model
runs (two of them CNN-3) hit the 20-epoch cap without a clear validation-loss plateau.

% =====================================================================
\appendix
\section{Appendix}
% =====================================================================

\subsection{Additional figures}
\begin{figure}[H]
\centering
\includegraphics[width=\textwidth,height=0.42\textheight,keepaspectratio]{SV_comparison_bars.png}
\par\medskip
\includegraphics[width=\textwidth,height=0.42\textheight,keepaspectratio]{GT_comparison_bars.png}
\caption{Per-dataset metric and time bars: \run{SV} (top), \run{GT} (bottom).}
\end{figure}

\begin{figure}[H]
\centering
\includegraphics[width=\textwidth,height=0.42\textheight,keepaspectratio]{SV-TF-CNN-5_curves.png}
\par\medskip
\includegraphics[width=\textwidth,height=0.42\textheight,keepaspectratio]{GT-PT-CNN-5_curves.png}
\caption{Training curves of the best image models: \run{SV-TF-CNN-5} (top), \run{GT-PT-CNN-5} (bottom).}
\end{figure}

\begin{figure}[H]
\centering
\includegraphics[width=0.75\textwidth,height=0.42\textheight,keepaspectratio]{SV-TF-CNN-5_filters.png}
\par\medskip
\includegraphics[width=0.75\textwidth,height=0.42\textheight,keepaspectratio]{SV-PT-CNN-5_filters.png}
\caption{First-layer filters: \run{SV-TF-CNN-5} (top), \run{SV-PT-CNN-5} (bottom).}
\end{figure}

\begin{figure}[H]
\centering
\includegraphics[width=0.85\textwidth,height=0.42\textheight,keepaspectratio]{DB_correlation.png}
\par\medskip
\includegraphics[width=0.75\textwidth,height=0.42\textheight,keepaspectratio]{DB_categorical_rates.png}
\caption{\run{DB}: correlation matrix (top) and diabetes share per category (bottom).}
\end{figure}

\subsection{How to run the project}

\textbf{1. Environment} (Windows, macOS or Linux; Python 3.10--3.12):
\begin{lstlisting}
python -m venv .venv
.venv\Scripts\activate          # Windows
source .venv/bin/activate       # macOS / Linux
pip install -r requirements.txt
\end{lstlisting}
Regarding GPU support: Apple Silicon picks up \code{tensorflow-metal} and \code{mps} automatically; with an NVIDIA GPU,
PyTorch uses \code{cuda} (on Windows the CUDA build must be installed first, see below), whereas TensorFlow can only
use an NVIDIA GPU under Linux/WSL2. Absent a GPU, both frameworks fall back to the CPU.
\begin{lstlisting}
pip install torch --index-url https://download.pytorch.org/whl/cu126
\end{lstlisting}

\textbf{2. Data.} Download from Kaggle and place the files as follows (see \code{dataset/README.md}):
\begin{table}[H]
\centering
\small
\begin{tabular}{ll}
\toprule
Folder & Files \\
\midrule
\code{dataset/sv/} & \code{train_32x32.mat}, \code{test_32x32.mat} \\
\code{dataset/gt/} & \code{train.p}, \code{valid.p}, \code{test.p}, \code{signname.csv} \\
\code{dataset/db/} & \code{train.csv} (competition \code{playground-series-s5e12}; join it first) \\
\bottomrule
\end{tabular}
\end{table}

\textbf{3. Run the notebooks.} Open \code{sv_svhn.ipynb}, \code{gt_gtsrb.ipynb} or \code{db_diabetes.ipynb} from
\code{notebook/} in Jupyter or VS Code and execute all cells, or alternatively run from \code{notebook/}:
\begin{lstlisting}
jupyter nbconvert --to notebook --execute --inplace db_diabetes.ipynb
\end{lstlisting}
Every notebook is fully self-contained. The first run generates \code{dataset/<code>/splits.npz} (the 70/15/15 split),
after which it trains all four models and writes out \code{notebook/results/<run_id>.json}, saves the trained models
under \code{notebook/models/}, and produces the figures in \code{report/images/}. It is best to start with
\code{db_diabetes}, since it runs the fastest.

\end{document}
